\documentclass[10pt]{article}
\usepackage[T1]{fontenc}
\usepackage[utf8]{inputenc}
\usepackage[french]{babel}
\usepackage{amsmath,amssymb,mathrsfs,mathpazo}
\usepackage{geometry}
\geometry{a4paper,margin=2cm}
\usepackage{graphicx,array,tabularx,tabularray,colortbl,multirow,multicol,booktabs,enumitem,pifont,ulem,fancybox,diagbox,cancel}
\usepackage{pstricks,pstricks-add,pst-plot,pst-tree,pst-node,pst-text,pst-func,pst-math,pst-bspline,pst-3dplot,pst-eucl,pst-solides3d}
\let\pscalculate\undefined
\usepackage{pst-bezier}
\usepackage{tikz,pgfplots,tkz-tab}
\usetikzlibrary{arrows,arrows.meta,patterns,calc,positioning,trees,decorations.markings,shapes.geometric,shapes.misc,shapes.arrows,intersections,decorations.pathmorphing,decorations.pathreplacing,backgrounds,shadows,babel}
\pgfplotsset{compat=1.18}
\usepgfplotslibrary{fillbetween}
\usepackage[np]{numprint}
\usepackage{listings,makecell,pas-tableur,twoopt}
\let\barre\undefined
\usepackage[pstricks]{tablvar}
\definecolor{darkgreen}{RGB}{0,100,0}
\providecommand{\red}{\color{red}}
\providecommand{\blue}{\color{blue}}
\providecommand{\green}{\color{green}}
\providecommand{\black}{\color{black}}
\providecommand{\white}{\color{white}}
\providecommand{\gray}{\color{gray}}
\providecommand{\cyan}{\color{cyan}}
\providecommand{\magenta}{\color{magenta}}
\providecommand{\yellow}{\color{yellow}}
\providecommand{\violet}{\color{violet}}
\providecommand{\orange}{\color{orange}}
\providecommand{\euro}{\texteuro}
\providecommand{\EUR}{\texteuro}
\providecommand{\ieme}{\textsuperscript{e}}
\providecommand{\degres}{\ensuremath{^\circ}}
\providecommand{\degree}{\ensuremath{^\circ}}
\providecommand{\vv}[1]{\overrightarrow{#1}}
\providecommand{\Coord}[1]{#1}
\providecommand{\decosix}{}
\providecommand{\starredbullet}{\ensuremath{\bullet}}
\pagestyle{empty}
\setlength{\parindent}{0pt}
\usepackage{adjustbox}\DeclareUnicodeCharacter{2236}{:}\DeclareUnicodeCharacter{2190}{\ensuremath{\leftarrow}}
\providecommand{\R}{\mathbb{R}}
\providecommand{\N}{\mathbb{N}}
\providecommand{\Z}{\mathbb{Z}}
\providecommand{\Q}{\mathbb{Q}}
\providecommand{\D}{\mathbb{D}}
\providecommand{\C}{\mathbb{C}}
\providecommand{\e}{{\mathrm{e}}}
\providecommand{\vect}[1]{\overrightarrow{#1}}
\providecommand{\vectt}[1]{\overrightarrow{\mathrm{#1}}}
\providecommand{\vvt}[1]{\overrightarrow{\mathrm{#1}}}
\providecommand{\np}[2][]{{#2~#1}}
\providecommand{\vv}[1]{{\overrightarrow{#1}}}
\renewcommand{\d}{{\mathrm{d}}}
\providecommand{\degres}{{^\circ}}
\providecommand{\degre}{{^\circ}}
\providecommand{\blue}{\color{blue}}
\providecommand{\red}{\color{red}}
\providecommand{\enskip}{\quad}
\providecommand{\llbracket}{\lbrack\!\lbrack}
\providecommand{\rrbracket}{\rbrack\!\rbrack}
\providecommand{\euro}{{\text{€}}}
\providecommand{\no}{{\text{n°}}}
\providecommand{\interval}[2]{{[#1;#2]}}
\providecommand{\intervFO}[2]{{[#1;#2[}}
\providecommand{\dotfill}{\ldots}

\providecommand{\ds}{\displaystyle}
\providecommand{\textding}[1]{#1}
\providecommand{\Oij}{\(\mathrm{O};\vec{\imath},\vec{\jmath}\)}
\providecommand{\Oijk}{\(\mathrm{O};\vec{\imath},\vec{\jmath},\vec{k}\)}
\providecommand{\Ouv}{\(\mathrm{O};\vec{u},\vec{v}\)}
\providecommand{\intervFF}[2]{\ensuremath{[#1;#2]}}
\renewcommand{\euro}{\ifmmode\text{\texteuro}\else\texteuro\fi}
\DeclareUnicodeCharacter{00B0}{\ifmmode^{\circ}\else\textdegree\fi}
\DeclareUnicodeCharacter{20AC}{\ifmmode\text{\texteuro}\else\texteuro\fi}
% Énoncés : Éducation nationale. Transcription et corrigés : APMEP, Denis Vergès et contributeurs. Découpage et classement : CoopMaths, Sébastien Lozano, Éric Elter et contributeurs. Adaptation de la présentation : Maths Explorer.
% CC BY-SA : https://coopmaths.fr/www/about/
% https://coopmaths.fr/alea/static/e3c/2020/tex/e3c_2020_00_sujet22_4.tex
\begin{document}
\noindent\textit{APMEP / CoopMaths. CC BY-SA. Présentation adaptée par Maths Explorer.}\par\medskip
\medskip

\textbf{Les résultats seront donnés sous forme de fractions irréductibles.}

Une enquête a été menée auprès de lycéens pour estimer la proportion de ceux qui ont déjà consommé du cannabis. Pour encourager les réponses sincères, on met en place le protocole suivant :

Chaque adolescent lance d’abord un dé équilibré à 6 faces et l’enquêteur qui va l’interroger ne connaît pas le résultat du lancer. À la question « Avez-vous déjà consommé du cannabis ? », l’adolescent doit répondre :

\begin{itemize}
\item [\textbullet] \og non \fg{} si le résultat du lancer est 5, qu’il ait ou non déjà consommé du cannabis ;
\item [\textbullet] \og oui \fg{} si le résultat du lancer est 6, qu’il ait ou non déjà consommé du cannabis ;
\item [\textbullet] \og oui \fg{} ou « non » dans les autres cas, mais de façon sincère.
\end{itemize}

On note :

\begin{itemize}
\item [\textbullet] $N$ : l’évènement l’adolescent a répondu \og non \fg{} ;
\item [\textbullet] $O$ : l’évènement l’adolescent a répondu \og oui \fg{} ;
\item [\textbullet] $C$ : l’évènement l’adolescent a déjà consommé effectivement du cannabis ;
\item [\textbullet] $\overline{C}$ : l’évènement l’adolescent n'a jamais consommé du cannabis.
\end{itemize} 

Sur les lycéens qui ont participé à cette enquête on constate que la probabilité qu’un adolescent ait répondu \og oui \fg{} est de $\frac{3}{5}$ , soit $p(0) = \frac{3}{5}$.

On veut déterminer la probabilité, notée $p$, qu’un adolescent ait déjà consommé du cannabis.

On a donc $p(C)= p$.

\medskip

\begin{enumerate}
\item Justifier que la probabilité qu’un adolescent ait répondu « oui » sachant qu’il n’a jamais consommé de cannabis est $\frac{1}{6}$.
\item On a représenté ci-dessous l’arbre de probabilités représentant la situation. Compléter cet arbre.

\begin{center}
\psset{nodesepA=0pt,nodesepB=3pt,treesep=2cm,levelsep=3cm}
\pstree[treemode=R]{\TR{}}
{\pstree{\TR{$C$~~}\taput{$p$}}
	{
	\TR{$O$}\taput{$\dots$}
	\TR{$N$}\tbput{$\frac{1}{6}$}
	}
\pstree{\TR{$\overline{C}$~~}\tbput{$\dots$}}
	{\TR{$O$}\taput{$\frac{1}{6}$}
	\TR{$N$}\tbput{$\dots$}
	}
}
\end{center}
\item \begin{enumerate}
\item Démontrer que la probabilité $p$ qu’un adolescent ait déjà consommé du cannabis vérifie l’équation :

\[\frac{2}{3}p+\frac{1}{6}=\frac{3}{5}.\]

\item En déduire la valeur de $p$.
\end{enumerate}
\item Sachant qu’un adolescent a répondu « non » pendant l’enquête, quelle est la probabilité qu’il n’ait jamais consommé de cannabis ?
\end{enumerate}

\end{document}
